% Part: lambda-calculus
% Chapter: introduction
% Section: basic-pr-lambda

\documentclass[../../../include/open-logic-section]{subfiles}

\begin{document}

\olfileid[fr]{lam}{int}{bas}
\olsection{Les fonctions récursives primitives de base sont \texorpdfstring{$\lambd$}{λ}-définissables}

\begin{lem}
Les fonctions $\Zero$, $\Succ$ et $\Proj{n}{i}$ sont
$\lambd$-définissables.
\end{lem}

\begin{proof}
La fonction nulle~$\Zero$ est unaire : pour tout~$u$,
elle renvoie le numéral de Church
$\num{0}=\lambd[x][\lambd[y][y]]$.
Elle est donc définie par le terme
\[
Z\ident\lambd[u][\num{0}].
\]

La fonction successeur~$\Succ$ est définie par le terme
\[
\fn{Succ}\ident\lambd[u][\lambd[x][\lambd[y][x(uxy)]]].
\]
Pour un numéral~$\num{n}$, l'application
$\fn{Succ}\,\num{n}$ se réduit à
$\lambd[x][\lambd[y][x(\num{n}xy)]]$.
Le sous-terme $\num{n}xy$ applique $x$ exactement $n$~fois à~$y$.
L'application extérieure de $x$ ajoute une itération.
On obtient ainsi le numéral~$\num{n+1}$ par réduction~$\beta$.

Les projections ne présentent pas de difficulté :
$\fn{Proj}^n_i(x_0,\dots,x_{n-1})=x_i$.
Autrement dit, selon nos conventions,
$\fn{Proj}^n_i$ est le lambda-terme
$\lambd[x_0][\dots\lambd[x_{n-1}][x_i]]$.
\end{proof}

On peut distinguer le numéral~$\num{0}$ du terme~$Z$
en les appliquant tous deux à~$\num{1}$.
Le terme $\num{0}\,\num{1}$ se réduit à $\lambd[y][y]$,
tandis que $Z\,\num{1}$ se réduit au numéral~$\num{0}$.

\end{document}
